%% 14장 병렬 계산과 결정성
\chapter{병렬 계산과 결정성}
\label{ch:parallel}

\abstract{Apple M4 Pro에는 성능 코어 열 개와 효율 코어 네 개가 있다. 이 코어들을 함께 쓰면 계산이 빨라지지만, 공짜는 아니다. 일을 나누고 스레드를 깨우는 데 시간이 들고, 여러 스레드가 같은 메모리를 건드리면 결과가 실행 순서에 따라 달라질 수 있으며, 캐시 줄 하나를 두고 코어들이 다투면 오히려 느려진다. 이 장에서는 먼저 병렬화의 이득에 한계를 긋는 Amdahl의 법칙과 작업의 단위를 정하는 문제를 보고, 스레드 풀이 일을 나누어 주는 방법을 살핀다. 이어서 여러 스레드가 메모리를 함께 쓸 때의 규칙인 원자 연산과 메모리 순서, 캐시 일관성과 거짓 공유를 다룬다. 마지막으로 이 장의 중심 주제인 결정성, 곧 스레드의 수나 실행 순서와 상관없이 같은 비트를 내는 병렬 리덕션과 병렬 스캔을 설계하는 방법을 본다.}

\section{병렬화의 이득과 비용}
\label{sec:parallel-cost}

\subsection{Amdahl의 법칙}

어떤 계산에서 병렬화할 수 있는 부분이 실행 시간의 비율 $f$를 차지하고 나머지 $1-f$는 순차적으로만 할 수 있다고 하자. 코어 $p$개로 병렬 부분을 완벽하게 나누면 실행 시간은 $(1-f)+f/p$배가 되므로, 속도 향상은
\[
S(p)=\frac1{(1-f)+f/p}<\frac1{1-f}
\]
이다. 이것이 \term{Amdahl의 법칙}{Amdahl's law}이다\cite{amdahl1967}. 코어를 아무리 늘려도 속도 향상은 순차 부분이 정하는 상한 $1/(1-f)$를 넘지 못한다. $f=0.9$이면 코어가 열 개일 때 $S=1/(0.1+0.09)\approx5.3$이고, 무한히 많아도 $10$이다. 한편 문제의 크기를 코어 수에 맞추어 키우면 병렬 부분의 비율이 커지므로 큰 문제에서는 더 많은 이득을 볼 수 있다는 관점도 있는데, 이것을 Gustafson의 법칙이라 한다\cite{gustafson1988}. 두 법칙은 모순되지 않으며, 같은 크기의 문제를 빨리 푸는지 같은 시간에 더 큰 문제를 푸는지의 차이일 뿐이다.

\subsection{작업의 크기}

현실에서는 순차 부분 말고도 병렬화 자체의 비용이 있다. 잠든 스레드를 깨우는 데 시간이 들고, 일을 나누어 주고 결과를 모으는 데에도 시간이 들며, 나눈 조각들이 서로 다른 시간에 끝나면 빨리 끝난 코어는 기다린다. 이 비용은 대략 상수이므로, 일의 양이 작으면 병렬화가 오히려 느리다. 그래서 병렬 라이브러리는 일의 양이 어떤 \term{문턱값}{threshold}을 넘을 때만 나누고, 나눌 때에도 조각 하나의 크기가 너무 작지 않게 한다. 조각이 너무 크면 코어 사이의 균형이 나빠지고 너무 작으면 나누는 비용이 커지므로, 적당한 크기는 측정으로 정한다.

\begin{instardust}[\code{csrc/parallel/pool.hpp}]
stardust의 스레드 풀은 이 비용을 측정해 두었다. 주석에 따르면 잠든 일꾼 하나를 깨우는 데 호출한 스레드의 시간이 약 $0.7\mu$s 들고, 깨어난 일꾼은 약 $6\mu$s 뒤에야 일을 시작한다. 그래서 열 개의 스레드가 이득을 보려면 일의 총량이 약 $20\mu$s는 되어야 하고, 이 값이 상수 \code{kWakeWorkNs}이다. 언제 나눌지는 영역마다 자기 커널로 재서 정하는데, 배열 영역의 문턱값은 원소 $2^{15}$개이고 데이터프레임 영역은 행 $2^{17}$개이다. 이 문턱값을 넘는 일은 스레드가 하나뿐이어도 Python의 전역 인터프리터 잠금(GIL)을 풀고 계산하므로, 다른 Python 스레드가 그동안 계속 돈다.
\end{instardust}

\section{스레드 풀}
\label{sec:thread-pool}

운영체제의 스레드를 만들고 없애는 데는 수십 마이크로초가 들므로, 병렬 계산을 할 때마다 스레드를 새로 만들면 짧은 일에서는 그 비용이 일보다 크다. 그래서 프로그램이 시작할 때 일꾼 스레드들을 만들어 두고, 일이 없을 때는 재워 두었다가 일이 오면 깨워서 쓰는 \term{스레드 풀}{thread pool}을 쓴다. 재우는 방법에도 선택이 있다. 일이 끝난 직후에는 곧 다음 일이 올 가능성이 높으므로 잠깐 동안은 바쁘게 기다리면서(\term{스핀}{spin}) 새 일을 확인하고, 일정 시간이 지나도 일이 없으면 운영체제의 조건 변수로 잠든다. 스핀은 깨우는 지연이 없지만 코어를 차지하고, 잠은 코어를 내주지만 깨우는 데 시간이 든다.

일을 나누어 주는 가장 단순한 방법은 조각 $n$개를 스레드 $p$개에 미리 $n/p$개씩 나누는 \term{정적 분배}{static scheduling}이다. 그런데 코어마다 속도가 다르거나(성능 코어와 효율 코어), 조각마다 일의 양이 다르거나, 어떤 코어가 다른 프로그램에 잠시 빼앗기면 정적 분배는 가장 느린 스레드를 모두가 기다리게 만든다. \term{동적 분배}{dynamic scheduling}는 공유된 계수기 하나를 두고 각 스레드가 일을 마칠 때마다 계수기에서 다음 조각의 번호를 가져가게 한다. 빠른 스레드가 자연스럽게 더 많은 조각을 맡으므로 균형이 저절로 맞는다. 계수기를 원자적으로 하나씩 늘리는 연산이 필요한데, 이것이 다음 절의 주제이다. 조각들이 다시 하위 조각을 만들어 내는 재귀적인 병렬 계산에서는 스레드마다 자기 대기열을 두고 일이 떨어진 스레드가 남의 대기열에서 훔쳐 오는 \term{작업 훔치기}{work stealing}가 표준적인 방법이다\cite{blumofe1999}.

\begin{instardust}[\code{csrc/parallel/pool.cpp}]
stardust의 풀은 동적 분배를 쓴다. 일꾼들과 호출한 스레드 자신이 모두 원자 계수기 하나에서 \code{fetch\_add(1)}로 다음 조각의 번호를 가져가므로, 조각과 스레드 사이에 고정된 대응이 없다. 일꾼은 일이 끝나면 약 $100\mu$s 동안 스핀하다가 잠드는데, 이 시간은 반복 횟수가 아니라 시계로 잰다. ARM의 \code{yield} 명령은 한두 사이클이면 끝나서 큰 반복 횟수도 순식간에 지나가기 때문이다. 일꾼의 수는 기본적으로 성능 코어의 수에 맞춘다. 주석은 그 이유를 효율 코어가 느려서가 아니라(동적 분배라 느린 코어는 조각을 덜 맡을 뿐이다), 스레드를 열네 개까지 늘리면 메모리 대역폭에 묶인 연산이 $3$--$12\%$ 느려졌기 때문이라고 밝힌다.
\end{instardust}

\section{원자 연산과 메모리 순서}
\label{sec:atomics}

\subsection{데이터 경쟁}

두 스레드가 같은 변수 \code{count}에 \code{count = count + 1}을 동시에 하면 어떻게 될까. 이 문장은 읽기, 더하기, 쓰기의 세 단계이므로, 두 스레드가 같은 값을 읽고 각자 $1$을 더해 쓰면 증가가 한 번 사라진다. 이처럼 여러 스레드가 동기화 없이 같은 메모리에 접근하고 그중 적어도 하나가 쓰기인 것을 \term{데이터 경쟁}{data race}이라 하며, C++ 표준은 데이터 경쟁이 있는 프로그램의 동작을 정의되지 않은 것으로 둔다\cite{iso14882-2020}. 2장에서 본 부호 있는 정수의 넘침과 같은 종류의 함정이다.

\term{원자 연산}{atomic operation}은 다른 스레드가 중간 상태를 볼 수 없는 하나의 연산이다. \code{fetch\_add}는 값을 읽고 더해 쓰는 일을 원자적으로 하고 이전 값을 돌려주며, \code{compare\_exchange}(CAS)는 현재 값이 기대한 값과 같을 때만 새 값을 쓰고 성공 여부를 알려 준다. CAS가 있으면 거의 모든 원자적 갱신을 ``읽고, 새 값을 계산하고, 그사이에 아무도 바꾸지 않았으면 쓰고, 바뀌었으면 다시 하기''의 반복으로 만들 수 있다.

\begin{proposition}
\label{prop:cas-min}
원자 변수 $m$의 초깃값이 $+\infty$이고, 여러 스레드가 각자의 값 $r$에 대해 다음을 한다고 하자. $c\gets m$을 읽고, $r<c$인 동안 CAS로 $m$을 $c$에서 $r$로 바꾸려 시도하되, 실패하면 그때의 $m$을 $c$로 다시 읽는다. 그러면 모든 스레드가 끝난 뒤 $m$은 제시된 값들의 최솟값이며, 이 결과는 실행 순서와 무관하다.
\end{proposition}

\begin{proof}
$m$을 바꾸는 것은 성공한 CAS뿐이고, CAS는 $m=c$이고 $r<c$일 때만 $m$을 $r$로 바꾸므로 $m$은 줄어들기만 한다. 값 $r$을 제시한 스레드가 끝나는 경우는 둘이다. 읽은 $c$가 $r$ 이하여서 반복을 멈추었다면 그 순간 $m\le r$이었고 $m$은 줄어들기만 하므로 끝까지 $m\le r$이다. CAS가 성공했다면 그 순간 $m=r$이 되었으므로 역시 끝까지 $m\le r$이다. 따라서 마지막의 $m$은 제시된 모든 값 이하이다. 한편 $m$의 값은 언제나 초깃값이거나 어떤 스레드가 제시한 값이므로, 값이 하나라도 제시되었다면 마지막 $m$은 제시된 값 가운데 하나이면서 모두 이하인 값, 곧 최솟값이다.
\end{proof}

\begin{instardust}[\code{csrc/parallel/pool.hpp}]
\code{FirstRow::offer}가 바로 \cref{prop:cas-min}의 반복이다. 여러 조각이 병렬로 데이터를 검사하다가 잘못된 행을 찾으면 각자 그 행 번호를 제시하고, 가장 작은 번호가 남는다. 그래서 어떤 스레드가 먼저 끝나든 오류 메시지는 늘 같은 행을 가리킨다. 오류 메시지 하나에도 결정성을 지키는 것인데, 사용자가 같은 입력으로 같은 오류를 다시 볼 수 있어야 하기 때문이다.
\end{instardust}

\subsection{메모리 순서}

원자 연산 하나가 원자적이라는 것만으로는 부족하다. 한 스레드가 데이터를 쓴 뒤 ``준비됨''이라는 깃발을 세우고, 다른 스레드가 깃발을 본 뒤 데이터를 읽는다고 하자. 직관적으로는 깃발을 보았으면 데이터도 보여야 하지만, 컴파일러와 프로세서는 성능을 위해 서로 다른 주소에 대한 읽기와 쓰기의 순서를 바꿀 수 있다. 특히 ARM은 x86보다 이런 재배치를 훨씬 많이 허용하는 \term{약한 메모리 모델}{weak memory model}이다. 그래서 깃발이 먼저 보이고 데이터는 아직 옛 값이 보이는 일이 실제로 일어날 수 있다.

C++ 메모리 모델은 이 문제를 원자 연산에 붙이는 순서 조건으로 다룬다\cite{boehm2008cpp}. 가장 약한 \code{relaxed}는 그 변수 자체에 대한 원자성만 보장하고 다른 메모리와의 순서는 보장하지 않는다. 쓰기에 \code{release}를, 읽기에 \code{acquire}를 붙이면, acquire 읽기가 release 쓰기의 값을 본 경우 release 이전의 모든 쓰기가 acquire 이후의 모든 읽기에 보인다. 위의 깃발 예에서 깃발 쓰기를 release로, 깃발 읽기를 acquire로 하면 데이터가 반드시 보인다. 가장 강한 \code{seq\_cst}는 모든 스레드가 모든 seq\_cst 연산을 하나의 같은 순서로 본다는 \term{순차 일관성}{sequential consistency}을 보장하는데\cite{lamport1979}, 그만큼 비싸다. 계수기처럼 값 자체만 중요하고 다른 데이터를 넘기지 않는 원자 변수에는 relaxed로 충분하다.

\begin{instardust}[\code{csrc/parallel/pool.cpp}]
풀이 새 일을 일꾼들에게 알리는 방법이 release와 acquire의 전형적인 예이다. 호출한 스레드는 일의 함수 포인터와 조각의 수를 보통의 변수에 쓴 뒤, 세대 번호를 담은 상태 단어를 release로 갱신한다. 스핀하는 일꾼은 상태 단어를 acquire로 읽다가 세대가 바뀐 것을 보면, release 이전에 쓰인 일의 내용을 안전하게 읽는다. 반면 \code{FirstRow}의 행 번호와 잠든 일꾼의 수처럼 다른 데이터를 넘기지 않는 값은 relaxed로 다룬다.
\end{instardust}

\section{캐시 일관성과 거짓 공유}
\label{sec:false-sharing}

코어마다 자기 L1 캐시가 있으므로, 한 코어가 캐시에 있는 값을 바꾸면 다른 코어의 캐시에 있는 같은 주소의 사본은 낡은 것이 된다. 프로세서는 \term{캐시 일관성 프로토콜}{cache coherence protocol}로 이 문제를 처리한다. 대표적인 MESI 계열의 프로토콜에서 캐시 줄은 수정됨, 독점, 공유, 무효의 네 상태 가운데 하나를 갖고, 한 코어가 줄에 쓰려면 먼저 다른 모든 코어의 사본을 무효로 만들어야 한다\cite{papamarcos1984}. 쓰기가 잦은 줄은 코어 사이를 오가며 그때마다 상당한 지연을 낳는다.

문제는 일관성이 변수가 아니라 캐시 줄 단위로 관리된다는 데 있다. 서로 다른 스레드가 서로 다른 변수에 쓰더라도 두 변수가 같은 캐시 줄에 있으면 줄이 코어 사이를 오간다. 논리적으로는 아무것도 공유하지 않는데 하드웨어 수준에서는 공유하는 것이다. 이것을 \term{거짓 공유}{false sharing}라 한다. M4 Pro에서 네 스레드가 각자 자기 계수기를 relaxed \code{fetch\_add}로 늘리는 측정을 해 보면, 네 계수기가 연속된 $8$바이트 자리에 붙어 있을 때 연산 하나에 약 $27$ns가 걸리고, 각 계수기를 $128$바이트 경계에 따로 두면 약 $1.75$ns가 걸린다. $15$배 차이이며, 앞의 경우는 네 스레드가 한 계수기를 진짜로 공유할 때의 약 $24$ns와 거의 같다.

\begin{instardust}[\code{csrc/parallel/pool.cpp}]
풀은 조각 번호를 나누어 주는 원자 계수기와 일의 상태 단어를 각각 \code{alignas(128)}로 서로 다른 캐시 줄에 둔다. M4 Pro의 캐시 줄이 $128$바이트이기 때문이다. 일꾼들이 계수기를 두드릴 때마다 줄이 오가는데, 그 줄에 스핀하는 일꾼들이 읽는 상태 단어까지 있으면 일과 관계없는 읽기까지 줄을 따라 오가게 된다. 설계 문서는 bool 배열처럼 원소가 작은 연산에서 조각을 $4096$바이트로 잡았더니 원자 계수기와 캐시 줄을 주고받는 비용이 실제 일보다 커서 열 스레드가 한 스레드보다 느렸다는 기록도 남겨 두었다.
\end{instardust}

\section{결정적인 병렬 계산}
\label{sec:deterministic}

\subsection{문제}

9장에서 보았듯이 부동소수점 덧셈은 결합법칙을 만족하지 않으므로, 합하는 순서가 바뀌면 결과의 마지막 비트가 바뀔 수 있다. 병렬 합에서 가장 흔한 방법은 스레드마다 부분합을 구하고, 끝나는 대로 원자 연산이나 잠금으로 전체 합에 더하는 것이다. 이 방법은 스레드의 수가 바뀌면 부분합의 경계가 바뀌고, 같은 스레드 수에서도 끝나는 순서에 따라 더하는 순서가 바뀐다. 그래서 같은 입력으로 두 번 실행해도 다른 답이 나올 수 있다. 결과가 조금 다른 것이 그 자체로 틀린 것은 아니지만, 디버깅과 테스트가 어려워지고, 결과를 문턱값과 비교하는 분기가 실행마다 다른 길로 갈 수 있으며, 과학적 계산의 재현성이 무너진다\cite{demmel2013reproducible}.

\subsection{원칙}

결정적인 병렬 계산의 원칙은 두 가지로 요약된다. 첫째, 일을 나누는 방법은 데이터의 크기와 배치로만 정하고 스레드의 수로 정하지 않는다. 길이 $n$인 합을 언제나 같은 크기의 조각으로 나누면, 스레드가 몇 개든 각 조각이 계산하는 것은 같다. 둘째, 조각들의 결과를 모으는 순서는 고정한다. 조각 $0$, $1$, $2$의 순서로, 또는 고정된 나무의 모양대로 모으면, 어느 조각이 먼저 끝났는지는 결과에 영향을 주지 않는다. 동적 분배로 조각을 어느 스레드가 맡든 상관없이, 결과는 조각 번호에 따라 정해진 자리에 쓰고 마지막에 정해진 순서로 합친다.

이 원칙을 지키면 스레드 수가 달라도 같은 비트가 나온다. 그런데 이것은 생각보다 강한 약속이어서 대가가 따를 때가 있다. 스레드가 하나일 때 더 빠른 직렬 알고리즘이 있어도 병렬 판과 같은 비트를 내려면 같은 조각 구조를 써야 하기 때문이다. 그래서 약속의 범위를 ``같은 기계, 같은 스레드 수에서 같은 비트''로 좁히는 선택도 합리적이다. 이 경우에도 실행 순서에 따라 결과가 바뀌는 것, 곧 원자 연산의 순서에 기대는 합은 허용하지 않는다.

정수와 같이 결합법칙을 만족하는 연산은 이런 걱정이 없다. 정수의 합, 최댓값과 최솟값, 논리곱과 논리합은 조각을 어떻게 나누고 어떤 순서로 합쳐도 같으므로, 10장에서 본 정확한 정수 합은 결정성을 공짜로 얻는다. 최댓값의 위치를 찾는 \code{argmax}처럼 동점이 있는 연산은 ``앞 조각을 우선한다''는 규칙을 정하면 결합적인 연산이 된다.

\begin{instardust}[\code{docs/design/array.md} §10.5]
stardust의 약속은 ``같은 기계, 같은 스레드 수이면 실행마다 같은 비트''이다. 스레드 수가 달라도 같은 비트라는 더 강한 성질은 2026년 11월부터 약속에서 뺐지만, 배열 영역은 지금도 그렇게 만들어져 있다. 설계 문서는 연산의 종류마다 나누는 방법과 결과가 같은 이유를 표로 정리해 둔다. 원소별 연산은 원소마다 따로 계산하므로 어떻게 나누어도 같고, 축 리덕션은 출력이 걸친 축으로만 나누어 출력 하나에 입력이 접히는 순서가 직렬과 같게 하며, 출력이 하나인 긴 부동소수점 합은 9장의 쌍별 합 나무의 잎으로 나누어 같은 나무 순서로 합친다. 정수 연산과 최댓값, 최솟값은 고정 조각을 조각 순서대로 합친다. 테스트는 스레드가 $1$, $2$, $3$, $7$, $10$개일 때 같은 바이트가 나오는지 확인한다.
\end{instardust}

\section{병렬 스캔}
\label{sec:parallel-scan}

누적합 $y_i=x_1+\cdots+x_i$는 각 출력이 앞의 모든 입력에 의존하므로 병렬화할 수 없어 보인다. 그러나 결합법칙만 있으면 병렬로 계산할 수 있으며, 이 \term{스캔}{scan, prefix sum}은 병렬 알고리즘의 기본 구성 요소이다\cite{hillis1986,blelloch1990}. 코어 수가 많지 않은 공유 메모리 기계에서는 블록을 쓰는 세 단계 방법이 실용적이다.

\begin{algorithm}[t]
\begin{algorithmic}[1]
\Require 입력 $x_1,\dots,x_n$, 블록 길이 $B$
\Ensure 누적합 $y_i=x_1+\cdots+x_i$
\ForAll{블록 $b=0,1,\dots$ \textbf{in parallel}} \Comment{1단계: 블록마다 따로 스캔}
  \State 블록 $b$ 안에서 $y$를 블록의 첫 원소부터 차례로 누적한다
\EndFor
\State $c_0\gets0$; \textbf{for} $b\gets1,2,\dots$: $c_b\gets c_{b-1}+(\text{블록 }b-1\text{의 마지막 }y)$ \Comment{2단계: 직렬}
\ForAll{블록 $b\ge1$ \textbf{in parallel}} \Comment{3단계: 앞 블록들의 합을 더한다}
  \State 블록 $b$의 모든 $y_i$에 $c_b$를 더한다
\EndFor
\end{algorithmic}
\caption{블록을 쓰는 병렬 스캔}
\label{alg:block-scan}
\end{algorithm}

\cref{alg:block-scan}의 일의 양은 덧셈 약 $2n$번으로 직렬 스캔의 두 배이지만, 1단계와 3단계는 블록끼리 독립이므로 코어 수만큼 빨라지고, 2단계는 블록의 수만큼만 직렬이다. 정수의 스캔은 결합법칙이 성립하므로 결과가 직렬과 같다. 부동소수점 스캔은 블록의 경계에서 덧셈의 순서가 바뀌므로 직렬과 마지막 비트가 다를 수 있지만, 블록의 길이를 $n$에만 의존하게 정하면 스레드 수와 무관하게 같은 비트가 나온다. 2장에서 본 CSV 파서가 따옴표 안팎을 가르는 데 쓰는 XOR 누적도 결합적인 연산의 스캔이며, $64$비트 단어 안에서는 한 번의 다항식 곱셈으로 끝난다.

\begin{instardust}[\code{docs/design/array.md} §6.7]
배열의 \code{cumulative\_sum}은 부동소수점 레인이 $8192$개 이상이면 길이 $4096$의 블록으로 \cref{alg:block-scan}을 한다. 1단계에서는 블록 넷을 한 반복에서 나란히 돌려 서로 독립인 의존 사슬 넷을 동시에 진행시키는데, 13장에서 본 지연 시간 감추기이다. 설계 문서에는 메모리에 관한 흥미로운 기록이 있다. 네 블록에 원소를 하나씩 번갈아 저장했더니, 새로 할당한 출력 배열의 첫 쓰기에서 네 페이지의 페이지 부재가 겹쳐 $5.7$배 느려졌고, 그래서 블록마다 여덟 개씩 모아 쓰도록 바꾸었다고 한다. 블록의 구조가 $n$에만 의존하므로 결과는 스레드 수와 무관하다.
\end{instardust}

\section{맺음}

병렬 계산의 어려움은 계산 자체보다 그 주변에 있다. 일을 얼마나 잘게 나눌지, 스레드를 언제 깨울지, 공유된 메모리를 어떤 순서 규칙으로 다룰지, 캐시 줄 하나를 두고 코어들이 다투지 않게 하려면 데이터를 어디에 둘지가 모두 성능을 좌우한다. 그리고 부동소수점 산술에서는 여기에 결정성이 더해진다. 일을 나누는 방법을 데이터에만 의존하게 하고 결과를 모으는 순서를 고정하면, 병렬 계산도 직렬 계산처럼 언제나 같은 답을 낸다. 다음 장에서는 다시 한 코어 안으로 돌아가, 정렬과 선택을 중심으로 알고리즘의 분석 방법을 본다.

공유 메모리 병렬 프로그래밍의 원리는 \cite{herlihy2020}이, 메모리 모델의 형식적인 정의는 \cite{boehm2008cpp}가 자세하다.

\chapterreferences
